\section{Register sharing and the actual tape resources}
\label{sec:state-reduction-new}

\subsection{A nine-value obstruction from an admitted theorem-valued root}

The earlier register-lifetime report exhibits nine unequal live values inside
the Cantor unpairing macro, but explicitly does not claim that its chosen
formula registers satisfy the proof-search invariant. We strengthen that
particular obstruction with a small, exactly replayed source execution from
an admitted clean root. This remains an obstruction to direct pointwise
register identification, not to general eight-register algorithms.

In logical order $(N,T,W,A,D,H,X,Y,L)$, begin at
\[
  r_0=(4,2,3,0,6,10,0,0,1648).
\]
The formula registers are legitimate: $2$ codes $\forall\top$, $10$ codes
$\forall\forall\top$, and $3=P(2,0)$ codes a one-element stack containing
$\forall\top$. The three control temporaries $A,X,Y$ are zero. Thus this is
a clean root satisfying the inherited theorem/stack invariant; no consistency
assumption is needed to show that its displayed formulas are theorems.

The exact program reaches
\[
  r_*=(4,2,3,5,6,10,1,7,8)
\]
after 13,517 primitive source instructions, with decoded counter
\texttt{1100011001}. The next instruction is $\mathrm{CB}(1)$. Every one of
the nine displayed values is positive and they are pairwise distinct.

The calculation has a short conceptual explanation. The nonempty command
pointer is first decremented, and
\[
  1647=P(5,51),\qquad P(a,b)=\frac{(a+b)(a+b+1)}2+b.
\]
The first field read therefore sets $A=5$ and leaves $L=51$. During the
second field read, the Cantor walk has performed 43 updates when its values
are $X=1$, $Y=7$, and $L=8$, since $P(1,7)=43$ and $43+8=51$.
The six outer registers are unchanged during that walk.

\begin{proposition}[Direct aliasing fails even on admitted roots]
There is no implementation of the unchanged source program on every admitted
clean root that, at every source-instruction boundary, represents all nine
logical registers by the exact values of at most eight physical counters
through a map from logical registers to physical counters. The map may vary
with the source control location.
\end{proposition}
\begin{proof}
Apply the claimed representation to the displayed execution and its final
source boundary. Any map from nine logical registers to at most eight counters
identifies two logical registers. Their values would then be equal. The nine
values in $r_*$ are pairwise distinct, giving a contradiction.
\end{proof}

The proposition requires that each logical value itself be stored in its
assigned counter. It does not exclude recomputing a value, packing several
values into one counter, representing a value by a constant or an expression,
omitting a value that a revised simulation does not need, or moving to a
different register program. It also does not claim that $r_0$ occurs on the
literal blank computation of the 246-state table. These qualifications are
material; the proposition is not a lower bound on Turing-machine states.

The artifact \texttt{safe\_lifetime\_witness.py} has two independent finite
executors. One reads the exact 6,459-node literal tree in \texttt{Prog254.v}
and uses binary-counter path semantics. The other reads the 219-node
structured \texttt{Dag254.v}, using an explicit continuation stack instead
of binary-counter arithmetic. They agree on the instruction count, final
register vector, and next operation. The JSON record retains hashes of both
upstream programs. A second control has target
$(4,0,3,5,6,2,1,7,8)$, a permutation of $0,\ldots,8$, from an equally safe
root with $T=0$ and $H=2$.

\subsection{Pair merges with the actual spare restriction}

The previously completed broad pair-merge experiment used two spare values
$[0,2]$ and header distance 40. The actual 246-state bootstrap yields the
stronger resources of 167 value-one spares and header distance 181; increments
append value-one spares and the distance never decreases. A reduction that
requires these resources can therefore fail the older broader test.

We reran the complete pair family with those actual resource values and the
8,385 retained instruction representatives from the inherited interval
certificate. All six already excluded guards remain excluded. For each
unordered pair of the 246 states, identify the pair and independently select
the read-zero and read-one rows from either original state, redirecting all
incoming references. There are
\[
  4\binom{246}{2}=120,540
\]
recipes. The family permits arithmetic, selector, and dispatcher states to
participate. It holds the original logical-to-physical register permutation
fixed.

The baseline passes all 8,385 finite fixtures, exercises all 492 transition
pairs, and needs at most 2,692 moves for one fixture. Of the modified recipes,
76 pass the initial coverage probes; none passes the full affected-fixture
screen. Each candidate fixture has a 20,000-move bound and a represented tape
window from $-64$ through $1983$. The completed search takes 9.74 seconds in
the recorded environment. Its recipes and fixtures differ from the earlier
heterogeneous-spare experiment, although the mutation family is the same.

These are finite diagnostics. The representatives satisfy the supplied
interval restrictions and the stronger tape resources, but the experiment
does not prove that every representative occurs from blank. A timeout or
tape-window exit rejects a recipe for this screen without proving inequivalence.
Nor does one fixture with 167 spares prove a contract for every larger spare
count. No candidate survived to require such a universal proof or a new blank
initialization.

\subsection{A conserved quantity and the forward-carry collision}

Let $h$ be the absolute position of the register-file header, let
$s=\sum_{i=0}^{8}r_i$ be the sum of the nine primary values, and let $k$ be
the number of spare blocks. At source instruction boundaries of the original
246-state representation,
\[
  h+s-k=15.
\]
The initial values are $h=176$, $s=6$, and $k=167$. An increment raises both
$s$ and $k$ by one. A positive decrement raises $h$ by one and lowers $s$ by
one. Zero decrement, break, continue, and the counter-overflow adjustment
preserve the quantity; counter overflow shifts the left marker, not the
absolute header. Induction on source steps proves the formula.

This identity cannot by itself repair the information loss of the previously
proposed 244-state forward-carry decrement. Its zero and positive collision
inputs have the same header and all other primary values, while the positive
input has both selected-register value and spare count larger by one.
Consequently both inputs have the same value of $h+s-k$.

Moreover, absolute header position is not directly observable from a pointed
runtime frame once the original blank origin is unmarked. The available
distance $d$ from the current dispatcher head to the header satisfies
\[
  d+s-k=20+o,
\]
where $o$ is the number of program-counter overflows since bootstrap.
Any repair that proposes to use spare count as a branch flag must specify a
stronger, accessible invariant and an actual procedure for recovering the
flag. A conserved affine expression alone supplies neither.
